%===========================================================================================\\
%============================ EP001_physics_system.tex ====================================\\
%===========================================================================================\\
% EP001: Physical System Description - Unified Learning Module
% Consolidated from: LT-7 Paper + E003 Slides + Thesis
%===========================================================================================\\

\documentclass[11pt,a4paper]{article}

% Packages (hyperref must be last)
\usepackage[utf8]{inputenc}
\usepackage[T1]{fontenc}
\usepackage{amsmath,amssymb,amsthm,amsfonts}
\usepackage{geometry}
\usepackage{graphicx}
\usepackage{booktabs}
\usepackage{array}
\usepackage{enumerate}
\usepackage{fancyhdr}
\usepackage{xcolor}
\usepackage{listings}
\usepackage{titlesec}
\usepackage{tikz}
\usetikzlibrary{arrows,calc,positioning}

% Page geometry
\geometry{
    left=2.5cm,
    right=2.5cm,
    top=2.5cm,
    bottom=2.5cm,
    headheight=15pt
}

% Color definitions
\definecolor{moduleblue}{RGB}{0, 102, 153}
\definecolor{highlightgreen}{RGB}{0, 128, 0}
\definecolor{noteblue}{RGB}{0, 153, 204}
\definecolor{lightblue}{RGB}{230, 240, 255}

% Hyperref setup (with compatibility)
\usepackage[hypertexnames=false, colorlinks=true, linkcolor=moduleblue, filecolor=magenta, urlcolor=moduleblue, pdftitle={EP001: Physical System Description}, pdfauthor={Double-Inverted Pendulum Learning Module}]{hyperref}

% Header/footer setup
\pagestyle{fancy}
\fancyhf{}
\fancyhead[L]{EP001: Physical System Description}
\fancyhead[R]{Learning Module}
\fancyfoot[C]{\thepage}
\renewcommand{\headrulewidth}{0.4pt}

% Section formatting
\titleformat{\section}{\normalfont\Large\bfseries\color{moduleblue}}{\thesection}{1em}{}
\titleformat{\subsection}{\normalfont\bfseries\color{moduleblue}}{\thesubsection}{1em}{}
\titleformat{\subsubsection}{\normalfont\normalsize\bfseries}{\thesubsubsection}{1em}{}

% Custom environments
\newtheorem{theorem}{Theorem}[section]
\newtheorem{definition}{Definition}[section]
\newtheorem{proposition}{Proposition}[section]
\newtheorem{lemma}{Lemma}[section]

\newenvironment{note}[1][Note]{%
  \par\noindent\textcolor{noteblue}{\textbf{#1:}}\quad
}{\par}

\newenvironment{example}[1][Example]{%
  \par\noindent\textcolor{highlightgreen}{\textbf{#1:}}\quad
}{\par}

\lstset{
    basicstyle=\ttfamily\small,
    breaklines=true,
    frame=single,
    numbers=left,
    numberstyle=\tiny,
    numbersep=5pt,
    backgroundcolor=\color{lightblue},
    showspaces=false,
    showstringspaces=false,
    showtabs=false,
    tabsize=2,
    captionpos=b
}

\begin{document}

% Title page
\title{\textbf{\Large EP001: Physical System Description}\\[1em]
\large Double-Inverted Pendulum (DIP) System Model\\[1em]
\large Unified Learning Module\\[2em]
\small Consolidated from LT-7 Paper + E003 Slides + Thesis\\[1em]
\textbf{\large Learning Module}\\[3em]
Episode ID: EP001 \\
Difficulty: Intermediate (Advanced Undergraduate) \\
Duration: 2 hours \\
Status: Complete
}

\author{Double-Inverted Pendulum Control Project\\
Academic Learning Series}

\date{April 2026\\
\small Version 1.0 | Compiled: \today}

\maketitle
\thispagestyle{empty}
\newpage

% Table of Contents
\tableofcontents
\newpage

%===========================================================================================\\
% SECTION 1: INTRODUCTION
%===========================================================================================\\

\section{Introduction}

\subsection{Overview}

The \textbf{double-inverted pendulum (DIP)} represents one of the most challenging classical control problems in modern control theory. This nonlinear, underactuated mechanical system serves as a benchmark for testing advanced control strategies including:

\begin{itemize}
    \item Sliding Mode Control (SMC)
    \item Model Predictive Control (MPC)
    \item Reinforcement Learning approaches
    \item Adaptive and robust control methods
\end{itemize}

\subsection{System Characteristics}

The DIP system exhibits several challenging properties:

\begin{enumerate}
    \item \textbf{Nonlinearity}: Coupled, nonlinear differential equations of motion
    \item \textbf{Underactuation}: 2 degrees of freedom, 1 control input
    \item \textbf{Unstable equilibrium}: Open-loop unstable at upright position
    \item \textbf{Non-minimum phase}: Zero dynamics exhibit unstable behavior
    \item \textbf{Strong coupling}: Significant nonlinear coupling between links
\end{enumerate}

\begin{note}
The DIP system is commonly used in academic research and education to demonstrate advanced control concepts and compare different control strategies.
\end{note}

%===========================================================================================\\
% SECTION 2: PHYSICAL SYSTEM DESCRIPTION
%===========================================================================================\\

\section{Physical System Description}

\subsection{System Configuration}

The double-inverted pendulum consists of:

\begin{description}
    \item[Cart:] Moves along a 1D horizontal track with $\pm 1$ meter travel limit
    \item[Pendulum 1:] Rigid link pivoting at cart position, free to rotate $360^\circ$ ($\pm\pi$ rad)
    \item[Pendulum 2:] Rigid link pivoting at end of pendulum 1, free to rotate $360^\circ$
    \item[Actuation:] Single horizontal force $u$ applied to cart (motor-driven)
    \item[Sensing:] Encoders measure cart position $x$ and angles $\theta_1, \theta_2$; velocities estimated via differentiation
\end{description}

\subsection{Physical Constraints}

\begin{itemize}
    \item \textbf{Mass distribution:} $m_0 > m_1 > m_2$ (cart heaviest, tip lightest)
    \item \textbf{Length ratio:} $L_1 > L_2$ (longer base link provides larger control authority)
    \item \textbf{Inertia moments:} $I_1 > I_2$ (proportional to $m \cdot L^2$)
\end{itemize}

%===========================================================================================\\
% SECTION 3: MATHEMATICAL MODEL
%===========================================================================================\\

\section{Mathematical Model}

\subsection{State Vector Definition}

The state vector for the DIP system is defined as:

\begin{equation}
    \mathbf{x} = [x, \theta_1, \theta_2, \dot{x}, \dot{\theta}_1, \dot{\theta}_2]^T \in \mathbb{R}^6
\end{equation}

where:
\begin{itemize}
    \item $x$ -- cart position (m)
    \item $\theta_1$ -- angle of first pendulum from upright (rad)
    \item $\theta_2$ -- angle of second pendulum from upright (rad)
    \item $\dot{x}, \dot{\theta}_1, \dot{\theta}_2$ -- corresponding velocities
\end{itemize}

\subsection{Equations of Motion}

The nonlinear dynamics are derived using the \textbf{Euler-Lagrange method}, yielding:

\begin{equation}
    \mathbf{M}(\mathbf{q})\ddot{\mathbf{q}} + \mathbf{C}(\mathbf{q}, \dot{\mathbf{q}})\dot{\mathbf{q}} + \mathbf{G}(\mathbf{q}) + \mathbf{F}_{\text{friction}}\dot{\mathbf{q}} = \mathbf{B}u + \mathbf{d}(t)
\end{equation}

where $\mathbf{q} = [x, \theta_1, \theta_2]^T$ (generalized coordinates).

\subsection{Inertia Matrix}

The inertia matrix $\mathbf{M}(\mathbf{q}) \in \mathbb{R}^{3 \times 3}$ is symmetric and positive definite:

\begin{equation}
    \mathbf{M} = \begin{bmatrix}
        M_{11} & M_{12} & M_{13} \\
        M_{21} & M_{22} & M_{23} \\
        M_{31} & M_{32} & M_{33}
    \end{bmatrix}
\end{equation}

with elements:
\begin{align}
    M_{11} &= m_0 + m_1 + m_2 \\
    M_{12} &= M_{21} = (m_1 r_1 + m_2 L_1)\cos\theta_1 + m_2 r_2 \cos\theta_2 \\
    M_{13} &= M_{31} = m_2 r_2 \cos\theta_2 \\
    M_{22} &= m_1 r_1^2 + m_2 L_1^2 + I_1 \\
    M_{23} &= M_{32} = m_2 L_1 r_2 \cos(\theta_1 - \theta_2) + I_2 \\
    M_{33} &= m_2 r_2^2 + I_2
\end{align}

where $r_i$ = distance to center of mass, $I_i$ = moment of inertia.

\subsection{Coriolis and Centrifugal Matrix}

The matrix $\mathbf{C}(\mathbf{q}, \dot{\mathbf{q}}) \in \mathbb{R}^{3 \times 3}$ captures velocity-dependent forces, including:
\begin{itemize}
    \item \textbf{Centrifugal terms} $\propto \dot{\theta}_i^2$
    \item \textbf{Coriolis terms} $\propto \dot{\theta}_i \dot{\theta}_j$
\end{itemize}

\subsection{Gravity and Friction Terms}

\begin{equation}
    \mathbf{G}(\mathbf{q}) = \begin{bmatrix}
        0 \\
        (m_1 r_1 + m_2 L_1)g \cos\theta_1 \\
        m_2 r_2 g \cos\theta_2
    \end{bmatrix}, \quad
    \mathbf{F}_{\text{friction}} = \begin{bmatrix}
        b_0 & 0 & 0 \\
        0 & b_1 & 0 \\
        0 & 0 & b_2
    \end{bmatrix}
\end{equation}

%===========================================================================================\\
% SECTION 4: NONLINEARITY ANALYSIS
%===========================================================================================\\

\section{Nonlinearity Analysis}

\subsection{Configuration-Dependent Inertia}

The inertia matrix exhibits \textbf{strong configuration dependence}:

\begin{itemize}
    \item $M_{12}$ varies by up to \textbf{40\%} as $\theta_1$ changes from 0 to $\pi/4$ (for $m_1=0.2$kg, $L_1=0.4$m)
    \item $M_{23}$ varies by up to \textbf{35\%} as $\theta_1-\theta_2$ changes
    \item This creates \textbf{state-dependent effective mass}
\end{itemize}

\begin{theorem}[Inertia Matrix Properties]
The inertia matrix $\mathbf{M}(\mathbf{q})$ for the DIP system satisfies:
\begin{enumerate}
    \item $\mathbf{M}(\mathbf{q}) = \mathbf{M}^T(\mathbf{q})$ (symmetric)
    \item $\mathbf{M}(\mathbf{q}) > 0$ (positive definite)
    \item $\dot{\mathbf{M}}(\mathbf{q}) - 2\mathbf{C}(\mathbf{q}, \dot{\mathbf{q}})$ is skew-symmetric
\end{enumerate}
\end{theorem}

\subsection{Trigonometric Nonlinearity}

Gravity terms exhibit \textbf{trigonometric nonlinearity}:
\begin{itemize}
    \item $\cos(\theta)$ terms in $\mathbf{G}(\mathbf{q})$
    \item Nonlinear coupling in Coriolis matrix
    \item Small-angle linearization invalid for large disturbances
\end{itemize}

\begin{example}[Small-Angle Linearization]
For small angles ($|\theta_i| < 0.1$ rad), the system can be approximated as:
\begin{equation}
    \mathbf{M}_0\ddot{\mathbf{q}} + \mathbf{C}_0\dot{\mathbf{q}} + \mathbf{G}_0\mathbf{q} = \mathbf{B}u
\end{equation}
where $\mathbf{M}_0 = \mathbf{M}(\mathbf{0})$, $\mathbf{G}_0 = \left.\frac{\partial \mathbf{G}}{\partial \mathbf{q}}\right|_{\mathbf{q}=0}$

\textbf{Warning:} Linearization fails for:
\begin{itemize}
    \item Large initial angles ($> \pm 0.3$ rad)
    \item Strong disturbances
    \item Swing-up maneuvers
\end{itemize}
\end{example}

%===========================================================================================\\
% SECTION 5: SYSTEM PARAMETERS
%===========================================================================================\\

\section{System Parameters}

\subsection{Nominal Parameter Values}

\begin{table}[h]
\centering
\begin{tabular}{c c l}
\toprule
\textbf{Parameter} & \textbf{Symbol} & \textbf{Value} \\
\midrule
Cart mass & $m_0$ & 1.0 kg \\
Pendulum 1 mass & $m_1$ & 0.2 kg \\
Pendulum 2 mass & $m_2$ & 0.1 kg \\
Pendulum 1 length & $L_1$ & 0.4 m \\
Pendulum 2 length & $L_2$ & 0.3 m \\
Center of mass 1 & $r_1$ & 0.2 m \\
Center of mass 2 & $r_2$ & 0.15 m \\
Inertia 1 & $I_1$ & 0.0133 kg·m$^2$ \\
Inertia 2 & $I_2$ & 0.00375 kg·m$^2$ \\
Cart friction & $b_0$ & 0.1 Ns/m \\
Joint friction 1 & $b_1$ & 0.05 Nms/rad \\
Joint friction 2 & $b_2$ & 0.03 Nms/rad \\
Gravity & $g$ & 9.81 m/s$^2$ \\
\bottomrule
\end{tabular}
\caption{Nominal DIP System Parameters}
\label{tab:nominal_parameters}
\end{table}

\subsection{Parameter Uncertainty}

Real systems exhibit parameter variations:
\begin{itemize}
    \item \textbf{Tolerance:} $\pm 10\%$ on masses and lengths
    \item \textbf{Friction:} $\pm 50\%$ uncertainty
    \item \textbf{Inertia:} $\pm 20\%$ uncertainty
\end{itemize}

\begin{note}
Robust control designs must account for these uncertainties to maintain stability and performance.
\end{note}

%===========================================================================================\\
% SECTION 6: LINEARIZED MODEL
%===========================================================================================\\

\section{Linearized Model}

\subsection{Equilibrium Points}

The DIP has \textbf{three equilibrium points}:

\begin{enumerate}
    \item \textbf{Upright-Upright:} $(\theta_1, \theta_2) = (0, 0)$ -- unstable, control objective
    \item \textbf{Upright-Down:} $(\theta_1, \theta_2) = (0, \pi)$ -- saddle point
    \item \textbf{Down-Down:} $(\theta_1, \theta_2) = (\pi, \pi)$ -- stable equilibrium
\end{enumerate}

\subsection{Linearized State-Space Representation}

Linearizing around the upright position $(\theta_1, \theta_2) = (0, 0)$:

\begin{equation}
    \dot{\mathbf{x}} = \mathbf{A}\mathbf{x} + \mathbf{B}u
\end{equation}

where:
\begin{align}
    \mathbf{A} &= \begin{bmatrix}
        \mathbf{0}_{3\times3} & \mathbf{I}_3 \\
        \mathbf{M}_0^{-1}\mathbf{G}_0 & \mathbf{M}_0^{-1}(\mathbf{C}_0 - \mathbf{F}_{\text{friction}})
    \end{bmatrix}, \\
    \mathbf{B} &= \begin{bmatrix} \mathbf{0}_{3\times1} \\ \mathbf{M}_0^{-1}\mathbf{B}_{\text{input}} \end{bmatrix}
\end{align}

\begin{proposition}[Controllability]
The linearized system $(\mathbf{A}, \mathbf{B})$ is \textbf{controllable}, meaning any state can be reached from any other state using appropriate control input.
\end{proposition}

%===========================================================================================\\
% SECTION 7: NUMERICAL PROPERTIES
%===========================================================================================\\

\section{Numerical Properties}

\subsection{State Scaling}

For numerical stability, states should be scaled:
\begin{itemize}
    \item Angles: radians (range $[-\pi, \pi]$)
    \item Positions: meters (range $[-1, 1]$)
    \item Velocities: rad/s or m/s (range $[-10, 10]$ typically)
\end{itemize}

\subsection{Time Scaling}

\begin{itemize}
    \item \textbf{Simulation step:} $\Delta t = 0.01$ s (100 Hz)
    \item \textbf{Control frequency:} 100--500 Hz
    \item \textbf{Real-time requirement:} Control computation $< 1$ ms
\end{itemize}

\begin{lemma}[Stiffness]
The DIP system exhibits \textbf{stiff behavior} near the upright position, requiring:
\begin{itemize}
    \item Small integration steps for accurate simulation
    \item Careful tuning of control gains
    \item Consideration of actuator bandwidth limits
\end{itemize}
\end{lemma}

%===========================================================================================\\
% SECTION 8: SUMMARY AND KEY TAKEAWAYS
%===========================================================================================\\

\section{Summary and Key Takeaways}

\subsection{System Overview}

The double-inverted pendulum is a:
\begin{itemize}
    \item \textbf{6-dimensional} state-space system
    \item \textbf{Nonlinear} dynamics with trigonometric terms
    \item \textbf{Underactuated} (1 input, 2 outputs)
    \item \textbf{Unstable} open-loop equilibrium
\end{itemize}

\subsection{Key Challenges}

\begin{enumerate}
    \item \textbf{Nonlinear coupling} between pendulum links
    \item \textbf{Parameter uncertainty} in mass and friction
    \item \textbf{Actuator constraints} (force limits, bandwidth)
    \item \textbf{Disturbance rejection} requirements
\end{enumerate}

\subsection{Control Implications}

\begin{itemize}
    \item Linear control methods require careful gain scheduling
    \item Nonlinear control (SMC, MPC) can handle larger operating ranges
    \item Robust designs must account for parameter variations
    \item Adaptive methods can compensate for uncertainties
\end{itemize}

\begin{theorem}[Control Design Trade-offs]
Control design for the DIP involves trade-offs between:
\begin{enumerate}
    \item Performance vs. robustness
    \item Complexity vs. implementability
    \item Linear vs. nonlinear approaches
    \item Model-based vs. data-driven methods
\end{enumerate}
\end{theorem}

%===========================================================================================\\
% END OF EP001
%===========================================================================================\\

\end{document}
